% ID: SEL-D0064-OA-E01
% ESTADO: validada. S064.
\textbf{a)} Para $0\le y\le5$, despejamos las fronteras: $(4-y)/3\le x\le6-y$. El cuadrilátero incluye sus lados. En la parte superior admite valores negativos de $x$.
\begin{center}
\setlength{\unitlength}{0.9mm}
\begin{picture}(96,61)(-22,-7)
\qbezier(-12,0)(28,0)(68,0)
\qbezier(0,0)(0,23)(0,46)
\qbezier(13.333,0)(36.667,0)(60,0)
\qbezier(60,0)(35,20)(10,40)
\qbezier(10,40)(3.333,40)(-3.333,40)
\qbezier(-3.333,40)(5,20)(13.333,0)
\put(69,0){$x$}\put(1,46){$y$}
\put(7,-6){\scriptsize $(4/3,0)$}
\put(57,-6){\scriptsize $(6,0)$}
\put(-20,40){\scriptsize $(-1/3,5)$}
\put(11,41){\scriptsize $(1,5)$}
\end{picture}
\end{center}
\textbf{b)} Las cotas de $x$ nunca se igualan en $[0,5]$. Para $y=0$ dan los vértices $(4/3,0)$ y $(6,0)$. Para $y=5$ dan $(-1/3,5)$ y $(1,5)$.

\textbf{c)} Evaluamos la función lineal en esos vértices, en el mismo orden: $20/3$, $30$, $40/3$ y $20$. El mínimo es $\boxed{20/3}$ sólo en $(4/3,0)$ y el máximo es $\boxed{30}$ sólo en $(6,0)$.
